% Part VI -- Algorithms, Optimization, and Machine Learning
% Chapter 20: Classification and model selection

\h{Classification and Model Selection}
\label{ch:20}

\booklink{ch:18}{Chapter 18} fitted one regression model and evaluated it once.
Real work is not one model. It is a search: several algorithms, several settings,
several thresholds, all competing for one recommendation. That search is an
optimization problem in the sense of \booklink{ch:19}{Chapter 19}, with one
difference that changes everything -- the score you are optimizing is measured on
data, so every comparison you run risks fitting the evaluation itself. This
chapter builds a classification workflow that survives that risk.

\begin{NoteBox}
\textbf{Prerequisites.} You should understand NumPy arrays, pandas DataFrames,
plotting, and the supervised-learning vocabulary of \booklink{ch:18}{Chapter 18}:
features $X$, target $y$, protected test data, baselines, and leakage. The
comfort-complaint dataset is reproducible synthetic data, used so the workflow
can be studied without privacy concerns.
\end{NoteBox}

\begin{tcolorbox}[title=Learning outcomes,colframe=orange,colback=white,arc=2mm]
After completing this chapter, you should be able to:
\begin{itemize}
    \item explain why accuracy misleads on an imbalanced target;
    \item build a majority-class baseline and read its confusion matrix;
    \item define precision, recall, and F1, and choose between them from the decision cost;
    \item assemble preprocessing and an estimator into one \c{Pipeline};
    \item separate a predicted probability from a predicted label;
    \item interpret ROC-AUC as ranking quality rather than accuracy;
    \item compare candidate models with \c{StratifiedKFold} cross-validation;
    \item read a validation curve to locate overfitting;
    \item tune hyperparameters with \c{GridSearchCV} without touching the test set;
    \item select a decision threshold from the relative cost of each error; and
    \item measure feature contribution with permutation importance and state its limits.
\end{itemize}
\end{tcolorbox}

\hh{Classification predicts a category, and categories are rarely balanced}

A classifier maps features to a category. The running example predicts whether a
room will generate an occupant comfort complaint in a given month, from
building and operating characteristics available before the complaint occurs.

The essential first measurement is not a model score. It is the share of rows in
each class.

\begin{lstlisting}[
    language=python,
    style=block,
    caption={Generate the reproducible comfort-complaint dataset},
    label={c20_data},
    captionpos=b
]
import numpy as np
import pandas as pd


def make_comfort_data(seed=42, sample_count=900):
    """Return synthetic features and a binary complaint target."""
    rng = np.random.default_rng(seed)
    features = pd.DataFrame({
        "floor_area_m2": rng.uniform(60, 520, sample_count),
        "glazing_ratio": rng.uniform(0.10, 0.68, sample_count),
        "occupant_density": rng.uniform(0.02, 0.16, sample_count),
        "setpoint_deviation_c": rng.normal(0.0, 2.4, sample_count),
        "outdoor_temp_c": rng.uniform(-8, 34, sample_count),
        "months_since_service": rng.integers(0, 36, sample_count),
    })
    risk = (
        -9.0
        + 1.10 * np.abs(features["setpoint_deviation_c"])
        + 4.5 * features["glazing_ratio"] * (features["outdoor_temp_c"] > 26)
        + 30.0 * features["occupant_density"]
        + 0.11 * features["months_since_service"]
        + rng.normal(0, 0.30, sample_count)
    )
    probability = 1 / (1 + np.exp(-risk))
    target = pd.Series(
        (rng.random(sample_count) < probability).astype(int),
        name="complaint",
    )
    return features, target


X, y = make_comfort_data()
print("Rows:", len(X))
print("Complaint rate:", round(y.mean(), 3))
\end{lstlisting}

Expected output:

\begin{lstlisting}[language=text, style=block, caption={An imbalanced target}, label={c20_data_output}, captionpos=b]
Rows: 900
Complaint rate: 0.239
\end{lstlisting}

\begin{SyntaxBox}{What the generator states explicitly}
\begin{itemize}
    \item \c{seed} makes the dataset reproducible, so every number in this
          chapter can be checked.
    \item \c{np.abs(setpoint_deviation_c)} means complaints rise when the room
          is too warm \emph{or} too cold. A straight line in that feature cannot
          represent this, which is why model choice will matter.
    \item The glazing term is multiplied by \c{(outdoor_temp_c > 26)}, so glazing
          only matters on hot days. That is an interaction.
    \item \c{rng.normal(0, 0.30, ...)} adds unexplained variation, so no model
          can reach perfect performance.
    \item The target is drawn from the probability rather than thresholded, so
          two identical rooms may differ. Real outcomes behave this way.
\end{itemize}
Because you wrote the rule, you know which features are real drivers
(\c{setpoint_deviation_c}, \c{occupant_density}, \c{months_since_service}, and
\c{glazing_ratio} on hot days) and which is pure noise (\c{floor_area_m2}).
Keep that answer key; you will check the model against it later.
\end{SyntaxBox}

\hh{Stratify the split so both classes survive it}

With 24 percent positives, an unlucky random split can leave the test set with a
noticeably different complaint rate than the training set. Stratifying preserves
the class proportions in both parts.

\begin{lstlisting}[
    language=python,
    style=block,
    caption={Split once, stratified by the target},
    label={c20_split},
    captionpos=b
]
from sklearn.model_selection import train_test_split

X_train, X_test, y_train, y_test = train_test_split(
    X,
    y,
    test_size=0.25,
    random_state=42,
    stratify=y,
)

print("Training rows:", len(y_train), "positives:", int(y_train.sum()))
print("Test rows:", len(y_test), "positives:", int(y_test.sum()))
\end{lstlisting}

Expected output:

\begin{lstlisting}[language=text, style=block, caption={Both parts keep the same complaint rate}, label={c20_split_output}, captionpos=b]
Training rows: 675 positives: 161
Test rows: 225 positives: 54
\end{lstlisting}

\c{stratify=y} is the one argument beginners most often omit. Without it, the
minority class count in the test set becomes another source of random variation,
and 54 positives is already a small number on which to base a claim.

\hh{Accuracy is the wrong first metric}

Before any model, measure what a rule with no information achieves. For
classification the honest baseline predicts the majority class every time.

\begin{lstlisting}[
    language=python,
    style=block,
    caption={The majority-class baseline and its confusion matrix},
    label={c20_baseline},
    captionpos=b
]
from sklearn.dummy import DummyClassifier
from sklearn.metrics import confusion_matrix

baseline = DummyClassifier(strategy="most_frequent")
baseline.fit(X_train, y_train)

print("Baseline accuracy:", round(baseline.score(X_test, y_test), 3))
print(confusion_matrix(y_test, baseline.predict(X_test)))
\end{lstlisting}

Expected output:

\begin{lstlisting}[language=text, style=block, caption={High accuracy, zero usefulness}, label={c20_baseline_output}, captionpos=b]
Baseline accuracy: 0.76
[[171   0]
 [ 54   0]]
\end{lstlisting}

\begin{ExplainBox}{Why 76 percent accuracy is worth nothing here}
The baseline predicts "no complaint" for all 225 test rooms. It is right about
all 171 rooms that had no complaint and wrong about all 54 that did, giving
$171/225 = 0.76$. Accuracy rewards it for ignoring the only cases anyone cares
about. Whenever one class dominates, accuracy is close to the majority share
regardless of skill, so a model must be judged on how it treats the minority
class. Report the baseline first, every time; it converts an impressive-sounding
score into a comparison.
\end{ExplainBox}

\hh{Read a confusion matrix before any single number}

The confusion matrix counts the four possible outcomes of a binary decision.
scikit-learn orders rows by actual label and columns by predicted label, both
ascending.

\begin{center}
\begin{tabular}{lll}
 & \textbf{Predicted 0} & \textbf{Predicted 1} \\
\textbf{Actual 0} & true negative & false positive \\
\textbf{Actual 1} & false negative & true positive \\
\end{tabular}
\end{center}

Name the two errors in the language of the problem before choosing a metric. A
false positive here means an unnecessary inspection visit: minor cost, wasted
technician time. A false negative means a room whose occupants complain with no
warning: the failure the system exists to prevent. Those two errors are not
equally bad, so no single symmetric score can rank models correctly on its own.

\hh{Precision, recall, and F1 answer different questions}

\begin{center}
\begin{tabular}{lll}
\textbf{Metric} & \textbf{Question it answers} & \textbf{Formula} \\
Precision & Of the rooms we flagged, how many were right? & $TP / (TP + FP)$ \\
Recall & Of the rooms that complained, how many did we find? & $TP / (TP + FN)$ \\
F1 & One number balancing the two & harmonic mean \\
\end{tabular}
\end{center}

Precision protects the technician's time. Recall protects the occupants. Raising
one usually lowers the other, so state which matters more in your application
before you compare models. F1 is a reasonable default only when you genuinely
have no preference.

\hh{Put preprocessing inside a pipeline}

Logistic regression is sensitive to feature scale: \c{floor_area_m2} spans
hundreds while \c{glazing_ratio} spans less than one. Standardizing fixes that,
but standardizing before the split leaks test information into the training
data, exactly the error described in \booklink{ch:18}{Chapter 18}. A
\c{Pipeline} solves both problems at once.

\begin{lstlisting}[
    language=python,
    style=block,
    caption={One object that scales and classifies},
    label={c20_pipeline},
    captionpos=b
]
from sklearn.linear_model import LogisticRegression
from sklearn.pipeline import Pipeline
from sklearn.preprocessing import StandardScaler

logistic = Pipeline([
    ("scale", StandardScaler()),
    ("model", LogisticRegression(max_iter=1000)),
])

logistic.fit(X_train, y_train)
print("Test accuracy:", round(logistic.score(X_test, y_test), 3))
\end{lstlisting}

\begin{SyntaxBox}{A pipeline is one estimator with several steps}
\begin{itemize}
    \item The argument is a list of \c{(name, step)} pairs. Every step except
          the last must transform; the last one predicts.
    \item \c{fit()} calls \c{fit_transform()} on each transformer using training
          data only, then fits the final estimator on the transformed result.
    \item \c{predict()} calls \c{transform()} with the parameters learned during
          fitting, so test rows are scaled using training means and standard
          deviations, never their own.
    \item Cross-validation refits the entire pipeline inside each fold, which is
          the only way to keep scaling honest during model selection.
\end{itemize}
Fitting a scaler on \c{X} and then splitting produces higher scores and a model
that will disappoint in deployment. The pipeline makes the correct order the
easy one.
\end{SyntaxBox}

\hh{Separate the predicted probability from the predicted label}

\c{predict()} returns labels. \c{predict_proba()} returns, for each row, the
estimated probability of each class. The label is produced from the probability
by comparing it with a threshold, and scikit-learn's default threshold is 0.5.

\begin{lstlisting}[
    language=python,
    style=block,
    caption={The label is a decision made from a probability},
    label={c20_proba},
    captionpos=b
]
probabilities = logistic.predict_proba(X_test)[:, 1]
labels = logistic.predict(X_test)

print("First three probabilities:", probabilities[:3].round(3))
print("Same rows as labels:", labels[:3])
print("Identical to a manual threshold:",
      np.array_equal(labels, (probabilities >= 0.5).astype(int)))
\end{lstlisting}

\c{[:, 1]} selects the second column, the probability of class 1, because
\c{predict_proba()} returns one column per class in the order given by
\c{model.classes_}. The final line prints \c{True}: the default classifier is
exactly the probability model plus a 0.5 cut. That cut is a business decision
you are allowed to change, and later in this chapter you will.

\hh{ROC-AUC measures ranking, not decisions}

Because the threshold is adjustable, a metric that judges the probabilities
themselves is useful. ROC-AUC is the probability that a randomly chosen
complaining room receives a higher score than a randomly chosen
non-complaining room. It is 1.0 for a perfect ranking, 0.5 for chance, and it
does not depend on the threshold.

\begin{lstlisting}[
    language=python,
    style=block,
    caption={Threshold-free and threshold-dependent metrics together},
    label={c20_metrics},
    captionpos=b
]
from sklearn.metrics import classification_report, roc_auc_score

print("ROC-AUC:", round(roc_auc_score(y_test, probabilities), 3))
print(classification_report(y_test, labels, digits=3))
\end{lstlisting}

For the logistic pipeline this reports ROC-AUC 0.776, accuracy 0.831, precision
0.750, and recall 0.444. Read those together: the model ranks rooms far better
than chance, and it is right three times in four when it does flag a room, but
at the default threshold it misses more than half of the complaints. Accuracy
0.831 against the 0.760 baseline barely registers that failure.

\begin{ExplainBox}{A good ranking and a bad decision are different problems}
ROC-AUC only asks whether complaining rooms score above non-complaining rooms.
It says nothing about where you cut the list. A model can have excellent ROC-AUC
and terrible recall at threshold 0.5, which is precisely what happens on an
imbalanced target: the estimated probabilities are informative but mostly below
0.5, so almost nothing crosses the default line. Diagnose ranking with ROC-AUC,
then set the threshold separately.
\end{ExplainBox}

\hh{Compare candidates with cross-validation, not with the test set}

You now want to know whether a different model would do better. Every answer you
get from the test set spends some of its independence: after ten comparisons, the
winner is partly the model that happened to suit those 225 rows.
Cross-validation answers the question inside the training data instead.

$k$-fold cross-validation splits the training rows into $k$ parts, fits on
$k-1$ of them, scores on the held-out part, and repeats so every row is scored
once. Stratified folds preserve the class balance in each part.

\bookfigure[0.98\textwidth]{Figures/Generated/ch20_classification_workflow.pdf}
{Candidate models and hyperparameters are compared with cross-validation inside the training data; the test set is opened once, at the end, to report the chosen model.}
{fig:ch20-workflow}

\begin{lstlisting}[
    language=python,
    style=block,
    caption={Score three candidates on the same folds},
    label={c20_crossval},
    captionpos=b
]
from sklearn.ensemble import RandomForestClassifier
from sklearn.model_selection import StratifiedKFold, cross_val_score
from sklearn.tree import DecisionTreeClassifier

folds = StratifiedKFold(n_splits=5, shuffle=True, random_state=42)

candidates = {
    "logistic pipeline": logistic,
    "decision tree (depth 5)": DecisionTreeClassifier(max_depth=5, random_state=42),
    "random forest": RandomForestClassifier(n_estimators=300, random_state=42),
}

for name, estimator in candidates.items():
    scores = cross_val_score(estimator, X_train, y_train, cv=folds, scoring="roc_auc")
    print(f"{name:<24} ROC-AUC {scores.mean():.3f} +/- {scores.std():.3f}")
\end{lstlisting}

Expected output:

\begin{lstlisting}[language=text, style=block, caption={Cross-validated comparison on training data only}, label={c20_crossval_output}, captionpos=b]
logistic pipeline        ROC-AUC 0.766 +/- 0.071
decision tree (depth 5)  ROC-AUC 0.741 +/- 0.034
random forest            ROC-AUC 0.836 +/- 0.041
\end{lstlisting}

\begin{SyntaxBox}{The spread is part of the result}
\c{cross_val_score()} returns one score per fold, so \c{scores.std()} measures how
much the estimate moves when the data change slightly. The logistic pipeline
varies by 0.071 across folds -- its individual fold scores range from 0.692 to
0.882. A gap between two models that is smaller than this spread is not evidence
of a difference. The forest's advantage of 0.07 over the logistic pipeline is
about one standard deviation, which is suggestive rather than decisive; more
folds or repeated cross-validation would sharpen it.
\end{SyntaxBox}

The forest wins because the target depends on the \emph{absolute} setpoint
deviation and on an interaction that only applies above 26 degrees. Trees split
on ranges, so they can represent both. Logistic regression fits one coefficient
per feature and cannot, which is why its ranking is weaker despite correct
scaling.

\hh{Read a validation curve to locate overfitting}

A hyperparameter that controls model flexibility -- here the maximum tree depth
-- can be swept while watching training and validation scores separately.

\bookfigure[0.98\textwidth]{Figures/Generated/ch20_model_selection.pdf}
{Training ROC-AUC rises to 1.000 as the tree deepens while cross-validated ROC-AUC peaks at depth 5 and then falls, and the tuned random forest ranks test rooms better than the logistic pipeline at every operating point.}
{fig:ch20-selection}

\begin{lstlisting}[
    language=python,
    style=block,
    caption={Sweep one hyperparameter and watch both curves},
    label={c20_validation_curve},
    captionpos=b
]
import numpy as np
from sklearn.model_selection import validation_curve

depths = np.arange(1, 16)
train_scores, valid_scores = validation_curve(
    DecisionTreeClassifier(random_state=42),
    X_train,
    y_train,
    param_name="max_depth",
    param_range=depths,
    cv=folds,
    scoring="roc_auc",
)

for depth, train_mean, valid_mean in zip(
    depths, train_scores.mean(axis=1), valid_scores.mean(axis=1)
):
    print(f"depth {depth:>2}  training {train_mean:.3f}  validation {valid_mean:.3f}")
\end{lstlisting}

Selected rows of the output:

\begin{center}
\begin{tabular}{lll}
\textbf{Maximum depth} & \textbf{Training ROC-AUC} & \textbf{Validation ROC-AUC} \\
1 & 0.651 & 0.574 \\
3 & 0.809 & 0.705 \\
5 & 0.905 & 0.741 \\
9 & 0.992 & 0.620 \\
15 & 1.000 & 0.611 \\
\end{tabular}
\end{center}

\begin{ExplainBox}{Both ends of the curve are failures}
At depth 1 the tree can ask one question, so it fits the training data poorly
and generalizes poorly. That is \textbf{underfitting}. At depth 15 it separates
the training rows perfectly -- ROC-AUC 1.000 -- by carving out regions that
contain a handful of rows each, and validation performance collapses to 0.611.
That is \textbf{overfitting}, and the widening gap between the two curves is its
signature. The useful setting is the peak of the validation curve, depth 5, not
the peak of the training curve. A model evaluated only on the data it was fitted
to will always recommend the most complex option available.
\end{ExplainBox}

\hh{Ensembles average away part of the variance}

A single deep tree memorizes its training rows. A random forest fits many trees,
each on a bootstrap sample of the rows and choosing among a random subset of
features at each split, then averages their predicted probabilities. The
individual trees still overfit, but their errors are partly independent, so the
average is more stable. This is why the forest scored 0.836 where the tuned
single tree scored 0.741.

The cost is interpretability: you can read one depth-3 tree, but not 300 of them.
The permutation-importance measurement later in this chapter recovers part of
that explanation.

\hh{Tune hyperparameters on training data only}

\c{GridSearchCV} fits every combination of the requested settings with
cross-validation, keeps the best by mean validation score, and refits that
combination on the whole training set.

\begin{lstlisting}[
    language=python,
    style=block,
    caption={Search a small, justified grid},
    label={c20_gridsearch},
    captionpos=b
]
from sklearn.model_selection import GridSearchCV

search = GridSearchCV(
    RandomForestClassifier(random_state=42),
    param_grid={
        "n_estimators": [200, 400],
        "max_depth": [4, 8, None],
        "min_samples_leaf": [1, 5],
    },
    cv=folds,
    scoring="roc_auc",
    n_jobs=-1,
)
search.fit(X_train, y_train)

print("Best settings:", search.best_params_)
print("Best cross-validated ROC-AUC:", round(search.best_score_, 3))

tuned = search.best_estimator_
\end{lstlisting}

Expected output:

\begin{lstlisting}[language=text, style=block, caption={Twelve combinations, five folds, one winner}, label={c20_gridsearch_output}, captionpos=b]
Best settings: {'max_depth': 8, 'min_samples_leaf': 5, 'n_estimators': 400}
Best cross-validated ROC-AUC: 0.841
\end{lstlisting}

\begin{SyntaxBox}{What the search actually costs and returns}
\begin{itemize}
    \item The grid has $2 \times 3 \times 2 = 12$ combinations, and each is
          cross-validated over 5 folds, so 60 models are fitted.
    \item \c{scoring="roc_auc"} states the objective. Leaving it unset would
          optimize accuracy, which this chapter has already shown to be a poor
          target on an imbalanced problem.
    \item \c{n_jobs=-1} uses all available processor cores.
    \item \c{best_score_} is a \emph{cross-validated} score, not a test score.
    \item \c{best_estimator_} has already been refitted on all training rows, so
          it is ready to predict.
\end{itemize}
Grid size grows multiplicatively. Prefer a small grid you can justify, or
\c{RandomizedSearchCV} when the space is large.
\end{SyntaxBox}

\begin{ExplainBox}{Selection bias hides inside the best score}
\c{best_score_} is the maximum over twelve candidates, and a maximum over
candidates is optimistically biased: some of that 0.841 is the luck of matching
these particular folds. The test set exists to remove that bias, which is why the
test evaluation must come after tuning and must happen once. If you tune, look at
the test score, and then tune again, the test set has silently become a
validation set. Reporting an honest estimate of a tuning procedure requires
nested cross-validation, where the whole search runs inside an outer fold.
\end{ExplainBox}

\hh{Open the test set once}

\begin{lstlisting}[
    language=python,
    style=block,
    caption={The single final evaluation},
    label={c20_final_eval},
    captionpos=b
]
from sklearn.metrics import classification_report, confusion_matrix, roc_auc_score

tuned_probabilities = tuned.predict_proba(X_test)[:, 1]

print("Test ROC-AUC:", round(roc_auc_score(y_test, tuned_probabilities), 3))
print(confusion_matrix(y_test, tuned.predict(X_test)))
print(classification_report(y_test, tuned.predict(X_test), digits=3))
\end{lstlisting}

The tuned forest reaches test ROC-AUC 0.917 with accuracy 0.876, precision
0.861, and recall 0.574 at the default threshold. Its confusion matrix is 166
true negatives, 5 false positives, 23 false negatives, and 31 true positives.

Note that the test ROC-AUC of 0.917 is higher than the cross-validated 0.841.
That is not evidence that the model improved; the test rows simply turned out to
be easier than the training folds. With 54 positive test rows, the sampling
variation is large. Report both numbers and do not present the more flattering
one alone.

\hh{Choose the threshold from the cost of each error}
\label{sec:c20-threshold}

Recall 0.574 means the model misses 23 of 54 complaints. That is a property of
the threshold, not of the model, and the threshold is yours to set.

\bookfigure[0.98\textwidth]{Figures/Generated/ch20_threshold_tradeoff.pdf}
{At threshold 0.50 the tuned forest flags 36 rooms and finds 31 of 54 complaints; lowering the threshold trades precision for recall along a curve whose F1 peaks near 0.38.}
{fig:ch20-threshold}

\begin{lstlisting}[
    language=python,
    style=block,
    caption={Score several thresholds against the same probabilities},
    label={c20_threshold},
    captionpos=b
]
from sklearn.metrics import precision_recall_fscore_support

for threshold in (0.2, 0.3, 0.4, 0.5, 0.6):
    predicted = (tuned_probabilities >= threshold).astype(int)
    precision, recall, f1, _ = precision_recall_fscore_support(
        y_test,
        predicted,
        average="binary",
        zero_division=0,
    )
    print(f"threshold {threshold:.1f}  precision {precision:.3f}  "
          f"recall {recall:.3f}  F1 {f1:.3f}")
\end{lstlisting}

Expected output:

\begin{lstlisting}[language=text, style=block, caption={One model, five classifiers}, label={c20_threshold_output}, captionpos=b]
threshold 0.2  precision 0.516  recall 0.907  F1 0.658
threshold 0.3  precision 0.632  recall 0.796  F1 0.705
threshold 0.4  precision 0.809  recall 0.704  F1 0.752
threshold 0.5  precision 0.861  recall 0.574  F1 0.689
threshold 0.6  precision 1.000  recall 0.222  F1 0.364
\end{lstlisting}

Threshold 0.4 finds 38 of the 54 complaints instead of 31, at the price of some
additional unnecessary inspections. Whether that trade is right depends on the
cost of a visit against the cost of an unresolved complaint -- a question for the
facilities team, not the classifier.

\begin{ExplainBox}{Select the threshold without spending the test set}
The table above was computed on test data for teaching, so you could see the
mechanism. In real work, choose the threshold on validation folds or on a
separate calibration split, then apply it unchanged to the test set. Otherwise
the reported precision and recall describe a threshold that was chosen to look
good on exactly those rows. The same discipline that governs model selection
governs threshold selection.
\end{ExplainBox}

\hh{Ask which features the model actually used}
\label{sec:c20-importance}

Permutation importance shuffles one feature's values, breaking its relationship
with the target, and measures how much the score falls. A large drop means the
model relied on that feature.

\begin{lstlisting}[
    language=python,
    style=block,
    caption={Measure reliance by breaking one feature at a time},
    label={c20_importance},
    captionpos=b
]
from sklearn.inspection import permutation_importance

importance = permutation_importance(
    tuned,
    X_test,
    y_test,
    n_repeats=20,
    random_state=42,
    scoring="roc_auc",
)

for index in importance.importances_mean.argsort()[::-1]:
    print(f"{X.columns[index]:<22} {importance.importances_mean[index]:+.4f} "
          f"+/- {importance.importances_std[index]:.4f}")
\end{lstlisting}

Expected output:

\begin{lstlisting}[language=text, style=block, caption={Measured reliance, ranked}, label={c20_importance_output}, captionpos=b]
setpoint_deviation_c   +0.1492 +/- 0.0177
occupant_density       +0.1282 +/- 0.0262
months_since_service   +0.0938 +/- 0.0171
outdoor_temp_c         +0.0149 +/- 0.0062
glazing_ratio          +0.0016 +/- 0.0034
floor_area_m2          -0.0033 +/- 0.0033
\end{lstlisting}

Compare this with the answer key from Listing~\ref{c20_data}. The three main
drivers are recovered in the right order, and \c{floor_area_m2}, which the
generator never used, scores essentially zero. A small negative value simply
means shuffling a useless feature happened to help slightly; it is noise, not a
finding.

\begin{ExplainBox}{Importance is about the model, not about the building}
This measurement answers "how much does this fitted model rely on this column?"
It does not answer "how much does this quantity affect comfort?" Three cautions
follow. Correlated features share importance and can both look weak. A feature
the model never learned to use will score zero even if it matters physically.
And a high score never establishes that changing the feature would change the
outcome -- that is a causal claim requiring a different study design.
\end{ExplainBox}

\hh{Applied example: comfort-complaint triage}
\label{sec:c20-application}

The runnable program \c{examples/c20_comfort_complaint_classification.py}
performs the complete workflow: it generates the data, splits it stratified,
measures the majority-class baseline, cross-validates three candidates, tunes
the forest, evaluates the test set once, sweeps the threshold, computes
permutation importance, and exports a diagnostic figure.

\begin{lstlisting}[
    language=python,
    style=block,
    caption={The workflow in the order it must be executed},
    label={c20_application},
    captionpos=b
]
# 1. Protect the test set before anything else touches the data.
X_train, X_test, y_train, y_test = train_test_split(
    X, y, test_size=0.25, random_state=42, stratify=y
)
folds = StratifiedKFold(n_splits=5, shuffle=True, random_state=42)

# 2. Establish the no-information baseline.
baseline = DummyClassifier(strategy="most_frequent").fit(X_train, y_train)
baseline_accuracy = baseline.score(X_test, y_test)

# 3. Compare candidates with cross-validation, inside training data.
for name, estimator in candidates.items():
    scores = cross_val_score(estimator, X_train, y_train, cv=folds, scoring="roc_auc")
    print(f"{name:<24} ROC-AUC {scores.mean():.3f} +/- {scores.std():.3f}")

# 4. Tune the leading candidate, still inside training data.
search = GridSearchCV(
    RandomForestClassifier(random_state=42),
    param_grid={
        "n_estimators": [200, 400],
        "max_depth": [4, 8, None],
        "min_samples_leaf": [1, 5],
    },
    cv=folds,
    scoring="roc_auc",
    n_jobs=-1,
).fit(X_train, y_train)
tuned = search.best_estimator_

# 5. Open the test set once and report against the baseline.
probabilities = tuned.predict_proba(X_test)[:, 1]
print(f"Baseline accuracy {baseline_accuracy:.3f}")
print(f"Cross-validated ROC-AUC {search.best_score_:.3f}")
print(f"Test ROC-AUC {roc_auc_score(y_test, probabilities):.3f}")
\end{lstlisting}

\begin{SyntaxBox}{The order is the method}
\begin{enumerate}
    \item Splitting first means no later step can consult test outcomes.
    \item The baseline is computed before any model, so the comparison cannot be
          quietly dropped when it is unflattering.
    \item Candidates and hyperparameters are both judged on the same folds,
          which makes the comparison fair.
    \item The tuned estimator is refitted on all training rows automatically.
    \item The test set is read once, and both the cross-validated and test scores
          are reported so the reader can see the gap.
\end{enumerate}
Rearranging these steps does not produce a worse model. It produces a model
whose reported performance cannot be believed.
\end{SyntaxBox}

\hh{Document intended use, limits, and fairness}

A classification model that triggers an action needs written scope alongside its
metrics:

\begin{itemize}
    \item the decision it supports, who acts on it, and what they do;
    \item the data source, period, sampling rule, exclusions, and class balance;
    \item features available at prediction time and the exact target definition;
    \item the split strategy, baseline, cross-validated and test scores;
    \item the operating threshold and the cost reasoning behind it;
    \item precision and recall for groups that matter, not only overall, since a
          model can be accurate in aggregate and unreliable for one building
          type, floor, or occupant population;
    \item known distribution shifts, such as a new control system or a changed
          complaint-reporting policy, that invalidate the model; and
    \item the monitoring, retraining, and human-review requirements.
\end{itemize}

Every number in this chapter comes from synthetic data. They demonstrate that
the code recovers a rule that was deliberately planted; they are not evidence
about any real building.

\hh{Common mistakes and useful diagnoses}

\begin{itemize}
    \item \textbf{Accuracy is reported without a baseline.} On this data a
          useless rule scores 0.760. Always show the comparison.
    \item \textbf{\c{stratify} is omitted.} The minority class count then varies
          randomly between the split parts.
    \item \textbf{Scaling happens before the split.} Put transformers in a
          \c{Pipeline} so they are refitted inside every fold.
    \item \textbf{The test set is consulted repeatedly.} Each look spends
          independence; use cross-validation for development.
    \item \textbf{Only the tuned score is reported.} \c{best_score_} is a maximum
          over candidates and is optimistically biased.
    \item \textbf{Threshold 0.5 is treated as fixed.} It is a default, not a
          decision; choose it from the cost of each error.
    \item \textbf{ROC-AUC is read as accuracy.} It measures ranking quality and
          is independent of the threshold.
    \item \textbf{Importance is read as causation.} It measures model reliance
          on a column, nothing more.
    \item \textbf{Fold-to-fold spread is ignored.} A difference smaller than the
          spread is not a result.
\end{itemize}

\hh{Practice ladder}

\hhh{Level 0 -- Identify}

In the comfort example, identify the samples, the six features, the target, the
class balance, the baseline, the three candidate models, and the two scores that
must both be reported.

\hhh{Level 1 -- Modify}

Change \c{scoring} in the cross-validation from \c{"roc_auc"} to \c{"accuracy"}
and rerun the three candidates. Report the new ranking and explain which
comparison you would trust for this problem.

\hhh{Level 2 -- Explain}

The tuned forest scores 0.841 in cross-validation and 0.917 on the test set.
Explain two distinct reasons the test score could be higher, and state what
additional evidence would let you distinguish them.

\hhh{Level 3 -- Build}

Write \c{summarize_classifier(model, X_test, y_test, threshold)} that returns a
dictionary containing precision, recall, F1, ROC-AUC, and the four confusion
matrix counts at the given threshold. Use it to reproduce the threshold table.

\hhh{Level 4 -- Challenge}

An unresolved complaint costs roughly eight times an unnecessary inspection
visit. Write a function that computes the expected cost at each threshold from
0.05 to 0.95 and returns the cheapest threshold. Select it using cross-validated
probabilities from the training data only, then apply that threshold once to the
test set and report the result. Explain why selecting it directly on the test
set would overstate the saving.

\hh{Practice solutions}

\textbf{Level 0:} Samples are synthetic room-months; the features are floor
area, glazing ratio, occupant density, setpoint deviation, outdoor temperature,
and months since service; the target is the binary complaint flag; the class
balance is 23.9 percent positive; the baseline is the majority-class rule at
0.760 accuracy; the candidates are the logistic pipeline, a depth-5 decision
tree, and a random forest; both the cross-validated ROC-AUC and the single test
ROC-AUC must be reported.

\textbf{Level 2:} First, the test set contains only 54 positive rows, so the
score has large sampling variation and this split may simply be easier than the
average training fold. Second, cross-validated models are each fitted on 80
percent of the training data while the final model is refitted on all of it, so
the final model is slightly better trained than any fold model. Repeated
cross-validation with several random splits would separate these: if the gap
persists across many splits it reflects the refitting, and if it varies widely it
reflects sampling.

\textbf{Level 3}

\begin{lstlisting}[
    language=python,
    style=block,
    caption={One reusable classification summary},
    label={c20_solution_summary},
    captionpos=b
]
from sklearn.metrics import (
    confusion_matrix,
    precision_recall_fscore_support,
    roc_auc_score,
)


def summarize_classifier(model, X_test, y_test, threshold=0.5):
    if not 0.0 < threshold < 1.0:
        raise ValueError("threshold must lie strictly between 0 and 1")

    probabilities = model.predict_proba(X_test)[:, 1]
    predicted = (probabilities >= threshold).astype(int)
    precision, recall, f1, _ = precision_recall_fscore_support(
        y_test,
        predicted,
        average="binary",
        zero_division=0,
    )
    true_negative, false_positive, false_negative, true_positive = (
        confusion_matrix(y_test, predicted, labels=[0, 1]).ravel()
    )
    return {
        "threshold": threshold,
        "precision": precision,
        "recall": recall,
        "f1": f1,
        "roc_auc": roc_auc_score(y_test, probabilities),
        "true_negative": int(true_negative),
        "false_positive": int(false_positive),
        "false_negative": int(false_negative),
        "true_positive": int(true_positive),
    }
\end{lstlisting}

\c{labels=[0, 1]} forces the two-by-two shape even when a threshold produces no
predictions of one class, which would otherwise make \c{ravel()} fail. ROC-AUC is
computed from the probabilities, so it is the same at every threshold -- which is
exactly the point of including it.

\textbf{Level 4:} Choosing the threshold on the test set optimizes a parameter
against the same 225 rows used to report the saving, so the reported cost is a
minimum over thresholds rather than an estimate of future cost. Selecting it from
cross-validated training probabilities and then applying it unchanged keeps the
test evaluation independent.

\hh{Chapter checkpoint}

\begin{enumerate}
    \item Why is accuracy misleading when one class holds 76 percent of the rows?
    \item What does \c{stratify=y} protect?
    \item What question does precision answer, and what question does recall answer?
    \item Why must a scaler live inside a \c{Pipeline} rather than run before the split?
    \item What is the relationship between \c{predict_proba()} and \c{predict()}?
    \item What does ROC-AUC measure, and what does it deliberately ignore?
    \item Why is cross-validation preferred to the test set for comparing models?
    \item What are the two signatures of overfitting on a validation curve?
    \item Why is \c{best_score_} an optimistic estimate?
    \item What does permutation importance measure, and what does it not establish?
\end{enumerate}

\hh{Checkpoint answers}

\begin{enumerate}
    \item A rule that always predicts the majority class already scores 0.760
          while finding none of the cases of interest.
    \item It keeps the class proportions of the full dataset in both the training
          and test parts.
    \item Precision: of the rows flagged, how many were correct. Recall: of the
          rows that were truly positive, how many were found.
    \item So the scaling parameters are learned from training rows only, and are
          refitted inside each cross-validation fold.
    \item \c{predict()} is \c{predict_proba()} compared against a threshold,
          which defaults to 0.5.
    \item It measures how well the model ranks positives above negatives, and it
          ignores where the decision threshold is placed.
    \item Because every comparison made on the test set reduces its independence,
          and the final score would then reflect selection rather than
          generalization.
    \item The training score keeps rising toward 1.0 while the validation score
          falls, and the gap between the curves widens.
    \item It is the maximum over many candidates, so part of it is the luck of
          matching those particular folds.
    \item How much a fitted model relies on a column; it does not establish
          physical importance or causation.
\end{enumerate}

\begin{tcolorbox}[title=Chapter summary,colframe=orange,colback=white,arc=2mm]
Classification is a decision built from a ranking, and the two must be judged
separately: use ROC-AUC and cross-validation to choose the model, then choose
the threshold from the cost of each error. Model selection is itself a search,
so it must run inside the training data; the test set is opened once, reported
beside the cross-validated score, and never used to pick anything.
\booklink{ch:21}{Chapter 21} adds one more candidate to this comparison -- a
neural network -- and holds it to exactly the same standard of evidence.
\end{tcolorbox}
